\section*{Bab \ref{ch:g12:age-of-the-earth} --- Penanggalan Radiometrik dan Umur Bumi}

\begin{solution}{exo:g12:age-of-the-earth:1}
(a) \omterm{def:g12:age-of-the-earth:closed-system}{sistem tertutup}; (b) \omterm{def:g12:age-of-the-earth:isochron}{isokron}; (c) \omterm{def:g12:age-of-the-earth:age}{umur Bumi}.
\end{solution}

\begin{solution}{exo:g12:age-of-the-earth:2}
Plagioklas (0.35), hornblende (1.20), batuan utuh (3.10), feldspar-K (6.40),
biotit (24.0). Biotit terplot paling tinggi: biotit mengandung rubidium-87
paling banyak untuk stronsium-86-nya, sehingga rubidiumnya telah membuat
stronsium-87 paling banyak sejak granit itu terbentuk.
\end{solution}

\begin{solution}{exo:g12:age-of-the-earth:3}
$T = \ln 2/\lambda = 0.6931/\num{1.3972e-11} \approx \num{4.96e10}$ tahun,
49.6 miliar tahun: $49.6/4.54 \approx 11$ kali \omterm{def:g12:age-of-the-earth:age}{umur Bumi}.
\end{solution}

\begin{solution}{exo:g12:age-of-the-earth:4}
$\lambda t = \num{1.3972e-11} \times \num{3.10e8} \approx 0.004331$, dan
$e^{0.004331} - 1 \approx 0.004341$. Untuk 480 juta tahun, $\lambda t
\approx 0.006707$ dan $e^{\lambda t} - 1 \approx 0.006729$. Keduanya dekat
dengan $\lambda t$ itu sendiri, karena $\lambda t$ kecil.
\end{solution}

\begin{solution}{exo:g12:age-of-the-earth:5}
Phillips (1860), Kelvin (1862), Joly (1899), Holmes (1927), Patterson
(1956). $4540/100 \approx 45$ kali.
\end{solution}

\begin{solution}{exo:g12:age-of-the-earth:6}
$0.7080 + 10.0 \times 0.004341 \approx 0.7514$ sekarang; 0.7080 ketika granit
itu terbentuk, seperti setiap mineral granit itu.
\end{solution}

\begin{solution}{exo:g12:age-of-the-earth:7}
(1) $m = (0.77525 - 0.71806)/(9.40 - 0.90) = 0.05719/8.50 \approx 0.006728$;
$t = \ln(1.006728)/\num{1.3972e-11} \approx \num{4.80e8}$ tahun, 480 juta
tahun. (2) $0.71806 + 1.70 \times 0.006728 \approx 0.72950$ dan $0.71806 + 4.20
\times 0.006728 \approx 0.74632$: keduanya terletak pada garis itu. (3) Ya:
granit memotong sekis, sehingga granit lebih muda, dan 310 kurang dari 480 juta
tahun.
\end{solution}

\begin{solution}{exo:g12:age-of-the-earth:8}
Semua titiknya akan jatuh di tempat yang sama: satu titik tidak menentukan
garis maupun kemiringan. Biotit terletak jauh di sebelah kanan yang lain;
galat kecil pada titik mana pun lalu jauh lebih sedikit mengubah kemiringan
garis yang melalui biotit dan plagioklas daripada kemiringan garis yang
melalui dua titik yang berdekatan.
\end{solution}

\begin{solution}{exo:g12:age-of-the-earth:9}
(1) Rubidiumnya tidak berubah, sehingga $x$-nya tetap; stronsium-87-nya
berkurang, sehingga $y$-nya turun: titik itu berpindah lurus ke bawah. (2) Garis
yang melalui titik itu dan plagioklas kurang curam, dan memberikan umur yang
terlalu muda. (3) Biotit tidak lagi terletak pada garis yang melalui
mineral-mineral lain: titik-titiknya tidak segaris, tanda bahwa satu sampel
tidak menjadi \omterm{def:g12:age-of-the-earth:closed-system}{sistem tertutup}.
\end{solution}

\begin{solution}{exo:g12:age-of-the-earth:10}
(1) $4540/96 \approx 47$ kali. (2) Bila sebagian garam meninggalkan laut,
misalnya terkubur di dalam sedimen, laut mengandung lebih sedikit garam daripada
seluruh garam yang dibawa sungai: membagi garam yang ada dengan pasokan
tahunannya memberikan waktu yang terlalu singkat.
\end{solution}

\begin{solution}{exo:g12:age-of-the-earth:11}
$m = (34.86 - 10.34)/(50.28 - 9.46) = 24.52/40.82 \approx 0.601$, di antara
kemiringan untuk 4.55 (0.590) dan 4.60 miliar tahun (0.610): sekitar 4.58
miliar tahun.
\end{solution}

\begin{solution}{exo:g12:age-of-the-earth:12}
(1) Basalt, dengan rasio yang lebih rendah, dari mantel; granit, dengan rasio
yang lebih tinggi, setidaknya sebagian dari kerak yang meleleh. (2) Rubidium
masuk ke dalam lelehan ketika mantel meleleh, sehingga rubidium terpusat di
dalam kerak, yang dari zaman ke zaman telah membuat lebih banyak stronsium-87.
(3) Tidak: mineral-mineral yang mengkristal dari satu magma mempunyai rasio
awal magma itu, dan kedua batuan ini tidak.
\end{solution}

\begin{solution}{exo:g12:age-of-the-earth:13}
(1) Rubidium-87: $2^{-4.54/49.61} \approx 0.94$, sekitar \qty{94}{\percent}
tersisa; uranium-238: $2^{-4.54/4.47} \approx 0.49$, sekitar separuhnya. (2) $m =
e^{\lambda t} - 1 \approx \num{1.3972e-11} \times 10^6 \approx
\num{1.4e-5}$: bahkan biotit granit Varden ($x = 24$) hanya akan naik
$0.0003$ di atas rasio awalnya. Stronsium-87 yang terbentuk begitu sedikit
sehingga kemiringannya sulit dibedakan dari nol: jam yang lambat cocok untuk
batuan tua.
\end{solution}

\begin{solution}{exo:g12:age-of-the-earth:14}
(1) $4.438 + 0.6026 \times 19.0 \approx 15.89$; $15.5$ lebih rendah $0.39$,
sekitar \qty{2.4}{\percent}: sedimen itu terletak dekat dengan garis. (2) 0.6026
terletak antara 0.590 (4.55) dan 0.610 (4.60): sekitar 4.58 miliar tahun, di
dalam rentang Patterson, 4.48 sampai 4.62 miliar tahun.
\end{solution}

\begin{solution}{exo:g12:age-of-the-earth:15}
(1) Simpanan itu lalu hanya kehilangan $46.7 - 24.4 = \qty{22.3}{TW}$ bersih,
sehingga simpanan yang sama bertahan $46.7/22.3 \approx 2.1$ kali lebih lama.
(2) $96 \times 2.1 \approx
200$ juta tahun, masih lebih dari dua puluh kali terlalu singkat. (3) \omterm{def:g12:internal-heat:convection}{Konveksi
mantel}, yang mengangkut panas ke atas dari bagian dalam dan menjaga gradien di
dekat permukaan tetap tinggi jauh lebih lama daripada konduksi saja.
\end{solution}

\begin{solution}{pb:g12:age-of-the-earth:1}
\textbf{1.} Biotit ($x = 24.0$).

\textbf{2.} Mineral-mineral itu mengkristal dari satu magma, dan mineral yang
mengkristal dari satu lelehan menyerap stronsiumnya dengan rasio isotopnya.

\textbf{3.} Stronsium-86 stabil dan tidak dihasilkan oleh peluruhan apa pun:
dalam \omterm{def:g12:age-of-the-earth:closed-system}{sistem tertutup} jumlahnya tidak berubah, sehingga stronsium-86 merupakan
tolok ukur yang tetap.

\textbf{4.} Stronsium-87, yang dihasilkan oleh peluruhan rubidium-87;
rubidium-87 telah berkurang dan stronsium-86 tetap sama.

\textbf{5.} Garis mendatar pada $y = 0.7080$: rasio yang sama, $x$ yang
berbeda-beda.

\textbf{6.} $(0.81218 - 0.70952)/(24.0 - 0.35) = 0.10266/23.65 \approx
0.004341$.

\textbf{7.} $(0.72146 - 0.70952)/(3.10 - 0.35) = 0.01194/2.75 \approx
0.004342$.

\textbf{8.} Keduanya sesuai sampai 1 banding 4000. Yang pertama lebih dapat
diandalkan: kedua titiknya berjauhan, sehingga pembulatan rasio-rasionya kurang
berpengaruh.

\textbf{9.} $0.70952 - 0.35 \times 0.004341 \approx 0.7080$.

\textbf{10.} $0.7080 + 6.40 \times 0.004341 \approx 0.73578$: tepat nilai
yang diukur.

\textbf{11.} $e^{\lambda t} = 1 + m = 1.004341$.

\textbf{12.} $\lambda t = \ln 1.004341 \approx 0.004331$.

\textbf{13.} $t = 0.004331/\num{1.3972e-11} \approx \num{3.10e8}$ tahun.

\textbf{14.} Ya: granit memotong sekis, sehingga menurut prinsip hubungan
potong-memotong granit lebih muda, dan 310 kurang dari 480 juta tahun.

\textbf{15.} 310 juta tahun terletak antara 358.9 dan 298.9 juta tahun:
zaman Karbon.

\textbf{16.} Rasio $^{87}\mathrm{Sr}/^{86}\mathrm{Sr}$ magma ketika granit itu
mengkristal.

\textbf{17.} Granit: $0.7080 > 0.7035$.

\textbf{18.} Rubidium masuk ke dalam lelehan ketika mantel meleleh, sehingga
rubidium terpusat di dalam kerak; dengan lebih banyak rubidium, kerak telah
membuat lebih banyak stronsium-87.

\textbf{19.} Bahwa magma itu berasal, setidaknya sebagian, dari \omterm{def:g12:continental-crust:continental-crust}{kerak benua} yang
meleleh, bukan hanya dari mantel.

\textbf{20.} Granit Varden berumur \textbf{310 juta tahun}.
\end{solution}
